Apply the finite additive Fourier idempotent for the prescribed
unit-squareclass subset modulo a sufficiently high power of $p$.
This idempotent keeps the $D_0$ coefficient. Its support condition
is stable under multiplication by unit squares, so it commutes
with the unit action and retains a scalar type, at a possibly
deeper level. At every prime outside $S_0$ take the spherical vector.
The packet there is unramified; its standard squarefree coefficient
$c_p^0$ is $1$ at both valuations zero and one
\cite[Section~I.3, pp.~377--378]{waldspurger}.
This includes the primes dividing $D_0$. Local Whittaker uniqueness
and the tensor-product factorization show that the resulting
global reference coefficient is still nonzero.

\paragraph{The exact full-space comparison.}
Let the resulting scalar character be $\psi$ and put
$\eta=\psi\chi_{-\sigma}$. Both have two-power order.
Let $V$ be the full good-Hecke eigenspace of this scalar character,
at the level just constructed. Keep all old multiplicities. In
Waldspurger's notation its integral-weight eigenform and character are
\[
 \phi_0=f_E\otimes\eta,\qquad
 \chi_0=\psi\chi_{-1},\qquad
 \phi_0\otimes\chi_0^{-1}=f_E\otimes\chi_\sigma .
\]
This also fixes the signs in the packet convention: Waldspurger's
positive index $D$ twists the last representation by $\chi_D$.
The Hecke eigenvalues in $V$ are $\eta(p)a_p(E)$.

The nonzero space implies hypothesis (H1) by Waldspurger's
Proposition~2 \cite[Section~I.2, p.~377]{waldspurger}. In explicit terms, at a principal-series
place its two inducing characters take value $1$ at $-1$.
Corollary~2 \cite[pp.~483--484]{waldspurger} applies to every $f\in V$ and any
positive squarefree pair whose ratio is a square in every
$\Q_p$ at a prime dividing the level. Our $D,D_0$ satisfy
this condition because the level has prime support $S_0$.
The final paragraph of that proof removes (H2) for this
same-$\Q_2$-squareclass comparison. No assumption of
non-supercuspidality at two is required.

\paragraph{The printed coefficient identity and its dyadic hypothesis.}
For clarity, write $\mathcal L_D=L(E^{(\sigma D)},1)$ and let
$N_{\rm W}$ denote the level in Waldspurger's notation, not the
conductor of $E$. In \cite[Corollary~2, pp.~483--484]{waldspurger}
take
\[
 k=3,\quad \chi=\psi,\quad \chi_0=\psi\chi_{-1},\quad
 \rho_0=\rho_E\otimes\eta,\quad n_1=D_0,\quad n_2=D.
\]
Then $\rho_0\otimes\chi_0^{-1}=\rho_E\otimes\chi_\sigma$.
The corollary, in its printed order, is
\[
 a_f(D_0)^2\mathcal L_D\,\psi(D/D_0)D^{1/2}
   =a_f(D)^2\mathcal L_{D_0}D_0^{1/2}.
\]
The passage from the automorphic central argument $1/2$ to the
classical argument $1$ is the same for both terms. Dividing only
by the nonzero reference value gives (E.10). In particular this
is a literal square identity, with neither complex conjugation
nor an unspecified absolute-value normalization.

Every prime of $N_{\rm W}$ belongs to $S_0$ at this point of the
construction. The prescribed local classes make $D/D_0$ a square
in $\mathbb Q_p^\times$ at each such prime, including $p=2$.
In the proof of the corollary, the local coefficient functions
at these primes have equal values at $D,D_0$. Outside the level,
the squarefree arguments have valuation zero or one and the
unramified coefficient function has value one in both cases.
The real coefficient contributes exactly the power
$(D/D_0)^{(k-2)/4}$ before squaring. The quotient of epsilon
factors is computed on pp.~483--484: its Hilbert-symbol product
disappears by reciprocity and its remaining character is
$\chi(D/D_0)=\psi(D/D_0)$. Thus these local terms do not leave
an additional prime-dependent factor in the displayed identity.

The final paragraph on p.~484 explicitly explains the use of
(H2): it is needed in that argument only to compare Fourier
indices in different squareclasses of $\mathbb Q_2^\times$.
Our pair is in one squareclass. Hence the corollary as stated
requires (H1) here, but does not impose (H2), a nonsupercuspidal
condition at two, or a new dyadic conductor restriction. The
nonzero scalar eigenspace supplies (H1) by
\cite[Proposition~2, p.~377]{waldspurger}. Moreover the proof on
p.~483 extends its linear relation to every vector in the full
eigenspace before squaring. The resulting identity therefore
applies simultaneously to old multiplicities, and later to both
Galois-conjugated and character-twisted vectors at a common level.
At that later level the same squareclass test is repeated; its
prime support is still $S_0$.
Putting $n_1=D_0,n_2=D$ in the corollary, with weight $3/2$, gives
\[
 a_f(D)^2=
 \psi(D/D_0)\,
 \frac{\sqrt D\,L(E^{(\sigma D)},1)}
      {\sqrt{D_0}\,L(E^{(\sigma D_0)},1)}
 a_f(D_0)^2.                                      \tag{E.10}
\]
The period cancels, so the quotient in this formula is the quotient
of the specified $L_*$-values. No product over the primes dividing
$D$ or $D_0$ is missing: the unramified factors at squarefree
valuations are exactly those checked above.

For clarity, (E.10) gives proportional coefficient functionals on
the whole space. If $\ell_0(f)=a_f(D_0)$ and $\ell_D(f)=a_f(D)$,
then $\ell_D^2=C\ell_0^2$ as a polynomial identity on $V$.
If $C=0$, then $\ell_D=0$. Otherwise, choose a square root $c$ of
$C$. The polynomial ring on the finite-dimensional vector space
$V$ is an integral domain, so
$(\ell_D-c\ell_0)(\ell_D+c\ell_0)=0$ implies
$\ell_D=c\ell_0$ or $\ell_D=-c\ell_0$.
The scalar is common to every vector. This argument does not
assume that $V$ is one-dimensional.
